\documentclass[11pt,a4paper]{article}
\usepackage[T1]{fontenc}
\usepackage[utf8]{inputenc}
\usepackage[brazil]{babel}
\usepackage{lmodern}
\usepackage[margin=2cm]{geometry}
\usepackage{microtype,booktabs,longtable,array,amsmath,acro,indentfirst,float}
\setlength{\parindent}{0.5cm}
\usepackage[hidelinks]{hyperref}
\DeclareAcronym{cpu}{short=CPU,long={unidade central de processamento, do inglês \textit{central processing unit}}}
\DeclareAcronym{gpu}{short=GPU,long={unidade de processamento gráfico, do inglês \textit{graphics processing unit}}}
\DeclareAcronym{ia}{short=IA,long=inteligência artificial}
\DeclareAcronym{sdr}{short=SDR,long={rádio definido por software, do inglês \textit{software-defined radio}}}
\newcolumntype{L}[1]{>{\raggedright\arraybackslash}p{#1}}
\setlength{\tabcolsep}{3pt}
\renewcommand{\arraystretch}{1.18}
\setlength{\emergencystretch}{2em}
\linespread{0.98}
\title{Impactos de desempenho em inteligência artificial e aplicações de rádio:\\ comparação entre ambientes com e sem Bitdefender}
\author{}
\date{}
\begin{document}
\maketitle
\section{Objetivo e limites da comparação}
Este relatório compara o desempenho de aplicações de \ac{ia} e operações de arquivos em dois computadores, um sem o Bitdefender (baseline) e outro com o produto em execução (Protected). As permissões disponíveis impediram testar o mesmo computador com a proteção ativada e desativada. Por isso, os resultados descrevem a diferença entre os ambientes e não isolam a contribuição exclusiva do Bitdefender. Os processadores, as unidades de armazenamento e os sistemas de arquivos também diferem, conforme a Tabela~\ref{tab:ambientes}. Mesmo com essa limitação, observa-se uma tendência de maior duração em várias cargas e operações de arquivos, acompanhada de consumo mensurável de \ac{cpu} pelos componentes do Bitdefender.

\begin{table}[H]\centering\small
\caption{Características dos computadores comparados.}\label{tab:ambientes}
\begin{tabular}{L{2.5cm}L{6cm}L{6cm}}\toprule & Baseline (sem) & Protected (com) \\ \midrule
Processador & Xeon Silver 4316, 20 núcleos / 40 threads & Xeon Gold 5220, 18 núcleos / 36 threads \\
Frequências & 2,30 GHz base / até 3,40 GHz turbo & 2,20 GHz base / até 3,90 GHz turbo \\
Memória & 64 GB DDR4, 2 módulos de 32 GB, 2666 MT/s & 64 GB DDR4, 4 módulos de 16 GB, 2666 MT/s \\
Unidade de dados & Samsung 980 PRO, 1 TB & LITEON CA5-8D512, 512 GB \\
Leitura sequencial & Até 7000 MB/s & Até 3300 MB/s \\
Gravação sequencial & Até 5000 MB/s & Até 2100 MB/s \\
Sistema operacional & Linux, Ubuntu & Linux, Ubuntu \\
\bottomrule\end{tabular}\end{table}
As frequências base e turbo são especificações da Intel para o \href{https://www.intel.com/content/www/us/en/products/sku/215270/intel-xeon-silver-4316-processor-30m-cache-2-30-ghz/specifications.html}{Xeon Silver 4316} e o \href{https://www.intel.com/content/www/us/en/products/sku/193388/intel-xeon-gold-5220-processor-24-75m-cache-2-20-ghz/specifications.html}{Xeon Gold 5220}.

\section{Bitdefender: recursos ocupados durante o trabalho}
A \ac{cpu} executa as instruções dos programas. Nas medições apresentadas, 100\% de \ac{cpu} corresponde a um processador lógico ocupado. Portanto, 200\% significa o equivalente a dois processadores lógicos.

Durante os intervalos monitorados da nova campanha, os processos visíveis do Bitdefender consumiram em média 196,30\% de \ac{cpu}, o equivalente a 1,96 processador lógico. O pico de um intervalo foi 617,46\%, equivalente a 6,17 processadores lógicos. Esses valores descrevem atividade real do agente durante o ensaio, inclusive preparação e cargas envolvendo a \ac{gpu}.

\begin{table}[H]\centering\small
\caption{Componentes do BitDefender. \ac{cpu} em porcentagem de um processador lógico. Memória residente média em MiB.}\label{tab:agente}
\vspace{0.1cm}
\begin{tabular}{L{6cm}rrr}\toprule Componente & \ac{cpu} média & \ac{cpu} pico & Memória \\ \midrule
Componente ctc & 142,41\% & 429,97\% & 1270,13 \\
Processo principal & 19,64\% & 73,98\% & 365,21 \\
Antimalware & 15,65\% & 246,00\% & 1255,45 \\
Componente networkmonitor & 13,28\% & 56,45\% & 125,36 \\
Componente fanotify & 5,10\% & 61,05\% & 59,61 \\
\bottomrule\end{tabular}\end{table}
O componente ctc concentrou 72,55\% do tempo de \ac{cpu} observado do agente. A Tabela~\ref{tab:agente} mostra que a demanda não se concentra apenas no componente antimalware. A soma da memória residente dos processos teve média de 3,20 GiB e máximo amostrado de 3,23 GiB. Essa soma pode contar páginas compartilhadas mais de uma vez e não mede memória privada exclusiva. Porém, ainda é um valor elevado.

Para uma aplicação que já ocupa quase toda a capacidade de processamento disponível, a demanda do agente reduz a folga para absorver variações. O pico também importa: uma média moderada pode coexistir com intervalos em que vários processadores lógicos atendem ao agente. A consequência depende dos núcleos compartilhados, das filas e dos prazos da aplicação.

\section{Resultados de desempenho}
Os resultados de treinamento apresentados na Tabela~\ref{tab:tempos} foram obtidos na \ac{cpu}. Nas cargas PyTorch em \ac{cpu} apresentadas, Protected levou mais tempo. A rede convolucional com dados gerados passou de 0,9666 s para 1,8065 s, aumento de 86,89\%. A taxa correspondente caiu de 6620,85 para 3542,71 amostras/s, redução de 46,49\%. Aumento de duração e redução de taxa são formas diferentes de expressar o mesmo trabalho fixo, com denominadores diferentes.

\begin{table}[H]\centering\small
\caption{Durações em segundos. Mediana $\pm$ desvio-padrão amostral, três repetições por condição. Variação positiva indica maior duração em Protected.}\label{tab:tempos}
\begin{tabular}{L{6.0cm}rrr}\toprule Tarefa & Baseline (sem) & Protected (com) & Variação \\ \midrule
\multicolumn{4}{l}{\textbf{\ac{cpu} e operações de arquivos}} \\
Inicialização do PyTorch & 0,7451 $\pm$ 0,0182 & 1,8406 $\pm$ 0,0526 & +147,04\% \\
Treino de rede convolucional / PyTorch & 0,9666 $\pm$ 0,0773 & 1,8065 $\pm$ 0,0580 & +86,89\% \\
Treino de linguagem / PyTorch & 13,8190 $\pm$ 0,4633 & 19,8194 $\pm$ 0,5534 & +43,42\% \\
Criação e leitura de 1024 arquivos & 0,0407 $\pm$ 0,0071 & 0,4816 $\pm$ 0,0192 & +1084,35\% \\
10000 leituras aleatórias & 0,0191 $\pm$ 0,0001 & 0,3091 $\pm$ 0,0716 & +1518,53\% \\
Gravação e leitura de controle, 64 MiB & 0,0536 $\pm$ 0,0054 & 0,3279 $\pm$ 0,0753 & +511,36\% \\
Gravação persistente de modelo / PyTorch & 0,0224 $\pm$ 0,0001 & 0,0307 $\pm$ 0,0049 & +37,25\% \\
Avaliação de 2000 exemplos / regressão logística & 0,0015 $\pm$ 0,0002 & 0,0054 $\pm$ 0,0024 & +260,25\% \\
\bottomrule\end{tabular}\end{table}

Nas operações de pequenos arquivos, o aumento foi de 40,67 ms para 481,62 ms por conjunto de 1024 gravações e 1024 leituras. São aproximadamente 441 ms adicionais por conjunto, não por arquivo. As 10000 leituras aleatórias passaram de 19,10 ms para 309,14 ms. Os controles usam arquivos recém-criados e podem ser atendidos parcialmente pelo cache, portanto não medem a velocidade física máxima das unidades.

Checkpoint de modelo é um arquivo que preserva o estado do modelo e do otimizador para restaurar o treinamento. Sua gravação persistente em PyTorch aumentou 37,25\%. Já a variante sem sincronização persistente ficou 25,81\% mais rápida em Protected. O controle de 64 MiB da tabela grava, sincroniza e lê um arquivo comum. Ele não serializa um modelo e não deve ser usado como estimativa universal do custo de checkpoints. Ou seja, um caso real pode ser mais pesado ainda. Além disso, esse tipo de operação é repetida centenas de vezes durante um treinamento real.

Cada valor reúne três repetições concluídas da mesma sessão. O desvio-padrão descreve a dispersão dessas repetições. As cargas de treinamento selecionadas usam 200 passos medidos e cinco de aquecimento. Foram comparadas tarefas com os mesmos parâmetros, conteúdo dos conjuntos de dados e versões das bibliotecas.

\section{O que essas ordens de grandeza significam em \texorpdfstring{\ac{ia}}{IA}}
\subsection{Treinamento: um pequeno custo pode se repetir milhões de vezes}
Treinar um modelo significa ajustar seus parâmetros usando sucessivos lotes de exemplos. Uma época corresponde a percorrer o conjunto de treinamento uma vez. Portanto, uma diferença pequena por lote ou por leitura pode se repetir muitas vezes. Milissegundos e meses se relacionam pela quantidade de repetições e pela parcela do tempo realmente afetada.

Os exemplos desta seção e da seção de rádio vêm do histórico real de testes e projetos do autor. Nesse histórico, houve treinamentos que chegaram a durar três meses. Um dos casos de treinamento envolve 14 milhões de amostras, organizadas em lotes de 32 exemplos. Cada época exige 437500 lotes. Um atraso adicional de 1 ms por lote representa 437,5 s, ou 7,29 minutos por época. Em 100 épocas, esse acréscimo soma 12,15 horas adicionais. O caso real mostra a escala de repetição das operações. O cálculo de 1 ms por lote quantifica o custo acumulado, sem atribuir esse atraso ao Bitdefender nos treinamentos anteriores.

Para quantificar o impacto em treinamentos que duram meses, o cálculo a seguir usa 30 dias como referência. Se uma etapa que ocupa todo esse tempo passasse a custar 86,89\% mais, como no caso PyTorch em \ac{cpu} da tabela, a projeção seria 56,07 dias. Se apenas 20\% do tempo original pertencesse à etapa afetada, a projeção cairia para 35,21 dias. Essas projeções quantificam o efeito acumulado de uma etapa mais lenta.

\subsection{Inferência: medir o modelo também exige medir o ambiente}
Inferência é usar um modelo já treinado para produzir uma resposta. Neste ensaio, a regressão logística apresentou 1,49 ms em baseline e 5,37 ms em Protected para avaliar um conjunto de 2000 exemplos. A rotina inclui a predição e o cálculo da acurácia. Esses tempos não são a latência de uma única inferência. A dispersão de 2,39 ms em Protected também mostra por que poucas medidas isoladas podem representar mal uma comparação.

Para um artigo, uma tarefa científica idêntica pode parecer mais lenta quando a medição inclui disputa com serviços de fundo. A consequência é uma comparação menos reprodutível ou uma conclusão sobre eficiência que o experimento não sustenta. Deve-se registrar hardware, carga concorrente e política de proteção, repetir as medidas e comparar modelos sob condições equivalentes. A diferença de ambiente pode ser confundida com uma diferença entre modelos. Em outras palavras, um artigo submetido em uma revista poderia ser rejeitado por causa desses resultados contaminados.

\section{Rádio: capacidade média e cumprimento de prazos}
Em \ac{sdr}, parte do processamento do sinal é executada pelo computador. As amostras chegam continuamente e precisam ser tratadas em tempo hábil. Aumentar a taxa de amostras ou a complexidade do processamento aproxima a aplicação do limite de \ac{cpu}. Nesse regime, o trabalho concorrente pode reduzir a taxa sustentável ou provocar acúmulo de dados e perdas, dependendo da margem dos buffers.

Em um rádio operando a 8 milhões de amostras por segundo, pequenas pausas já correspondem a grandes quantidades de dados. Nessa condição de operação, 1 ms corresponde a 8000 amostras recebidas e 500 ms correspondem a 4 milhões de amostras. Se uma pausa impedisse toda a cadeia de processar dados durante esse intervalo, seria necessário armazenar esse volume e depois recuperar o atraso. Nessa conversão, os números traduzem a duração de uma pausa na quantidade de amostras que precisaria aguardar processamento.

Outro caso do histórico de projetos é um estágio que recebe um bloco novo a cada 1 ms e leva 0,95 ms para processá-lo. Sua folga é de apenas 0,05 ms. Se o custo por bloco aumentar 10\%, o processamento passa a 1,045 ms, superior ao intervalo de chegada. A fila cresce enquanto essa condição persistir. Assim, uma diferença de apenas 0,095 ms por bloco pode comprometer a operação contínua mesmo que pareça desprezível em uma tarefa isolada.

Se a capacidade de processamento limita o enlace, reduzir o trabalho disponível por segundo pode exigir uma taxa de dados menor para manter o fluxo estável. Para calcular o efeito sobre a taxa de dados, considerando capacidade perfeitamente divisível, retirar 5,45\% da capacidade de uma aplicação já saturada transformaria um limite de 100 Mbit/s em cerca de 94,55 Mbit/s.

O pico observado, equivalente a 6,17 processadores lógicos, evidencia a importância da margem para variações. Sua duração e o escalonamento determinam se há aumento de latência, absorção pelos buffers ou perda de dados.

\section{Conclusão}
A comparação mantém uma tendência desfavorável em várias tarefas PyTorch e de arquivos. O dado mais direto sobre o Bitdefender é sua demanda concorrente: média de 1,96 processador lógico e pico amostrado de 6,17 durante a campanha. Para \ac{ia}, o efeito prático depende de quantas vezes a etapa se repete. Para rádio, depende de a execução terminar antes da chegada do próximo bloco.

\end{document}
